\section*{Chapitre \ref{ch:g9:ions} --- Les ions et les solutions ioniques}

\begin{solution}{exo:g9:ions:1}
Un \omterm{def:g7:atoms-and-molecules:atom}{atome} ou un groupe d'\omterm{def:g7:atoms-and-molecules:atom}{atomes} qui a perdu ou gagné des \omterm{def:g9:inside-the-atom:nucleus}{électrons} et qui
porte une charge. Un \omterm{def:g9:ions:ion}{cation} est positif (\omterm{def:g9:inside-the-atom:nucleus}{électrons} perdus), un \omterm{def:g9:ions:ion}{anion}
négatif (\omterm{def:g9:inside-the-atom:nucleus}{électrons} gagnés).
\end{solution}

\begin{solution}{exo:g9:ions:2}
\ce{Mg^{2+}}~: 12 \omterm{def:g9:inside-the-atom:proton}{protons}, 10 \omterm{def:g9:inside-the-atom:nucleus}{électrons}.
\end{solution}

\begin{solution}{exo:g9:ions:3}
\ce{KCl}, \ce{MgCl2}, \ce{CaO}.
\end{solution}

\begin{solution}{exo:g9:ions:4}
L'eau salée contient des \omterm{def:g9:ions:ion}{ions}, \ce{Na+(aq)} et \ce{Cl-(aq)}, qui se
déplacent et transportent des charges. Les \omterm{def:g7:atoms-and-molecules:molecule}{molécules} de sucre ne portent
pas de charge.
\end{solution}

\begin{solution}{exo:g9:ions:5}
\ce{CuSO4}, \ce{AlCl3}, \ce{Na2CO3}.
\end{solution}

\begin{solution}{exo:g9:ions:6}
Fer(III), \ce{Fe^{3+}}~: \ce{Fe^{3+}(aq) + 3OH-(aq) -> Fe(OH)3(s)}.
\end{solution}

\begin{solution}{exo:g9:ions:7}
Sulfate~: 2 \omterm{def:g9:inside-the-atom:nucleus}{électrons} en plus. Nitrate~: 1.
\end{solution}

\begin{solution}{exo:g9:ions:8}
Blanc avec le nitrate d'argent~: \ce{Cl-}. Bleu avec l'hydroxyde de
sodium~: \ce{Cu^{2+}}. C'est le chlorure de cuivre(II), \ce{CuCl2}.
\end{solution}

\begin{solution}{exo:g9:ions:9}
Quatre~: un de chaque côté (à gauche, à droite, au-dessus, au-dessous).
\end{solution}

\begin{solution}{exo:g9:ions:10}
\ce{Fe^{2+}(aq) + 2OH-(aq) -> Fe(OH)2(s)}.
\end{solution}

\begin{solution}{exo:g9:ions:11}
Le chlorure de zinc~: les \omterm{def:g9:ions:ion}{ions} sodium ne donnent pas de \omterm{def:g9:ions:precipitate}{précipité} avec
l'hydroxyde de sodium, les \omterm{def:g9:ions:ion}{ions} zinc en donnent un blanc~:
\ce{Zn^{2+}(aq) + 2OH-(aq) -> Zn(OH)2(s)}.
\end{solution}

\begin{solution}{exo:g9:ions:12}
\ce{CaCl2}~: 2000 \omterm{def:g9:ions:ion}{ions} chlorure. Charges~: $1000 \times (+2) + 2000 \times
(-1) = 0$.
\end{solution}

\begin{solution}{pb:g9:ions:1}
\textbf{1.} Fer(II), \ce{Fe^{2+}}~: \ce{Fe(OH)2}.

\textbf{2.} Fer(III), \ce{Fe^{3+}}~: \ce{Fe(OH)3}.

\textbf{3.} \ce{Fe^{2+}} en a un de plus~: $26 - 2 = 24$ \omterm{def:g9:inside-the-atom:nucleus}{électrons},
contre $26 - 3 = 23$ pour \ce{Fe^{3+}}.

\textbf{4.} \ce{Fe^{3+}(aq) + 3OH-(aq) -> Fe(OH)3(s)}.

\textbf{5.} Elle sort du puits avec des \omterm{def:g9:ions:ion}{ions} fer(II)~; au contact de
l'\omterm{def:g4:air-a-mixture-of-gases:air}{air}, ils deviennent des \omterm{def:g9:ions:ion}{ions} fer(III), qui forment le solide orange.

\textbf{6.} \ce{Fe(OH)3}~: $(+3) + 3 \times (-1) = 0$.

\textbf{7.} Pas à la sortie du puits~: les \omterm{def:g9:ions:ion}{ions} fer sont dissous et
traversent un filtre. Une fois le solide orange formé, oui~: c'est un
\omterm{def:g9:ions:precipitate}{précipité}, retenu par un filtre.

\textbf{8.} Non~: c'est un \omterm{def:g6:pure-substances-and-mixtures:mixture}{mélange} (de l'eau et des \omterm{def:g9:ions:ion}{ions} dissous).
Fraîchement tirée et limpide, elle est homogène.

\textbf{9.} $56 \times 1.67 \times 10^{-27} = \qty{9.35e-26}{kg}$.

\textbf{10.} $\qty{0.30}{mg} = \qty{0.30e-6}{kg} = \qty{3.0e-7}{kg}$.

\textbf{11.} Un \omterm{def:g9:ions:ion}{ion} diffère de l'\omterm{def:g7:atoms-and-molecules:atom}{atome} de deux ou trois \omterm{def:g9:inside-the-atom:nucleus}{électrons}, dont
la masse est négligeable devant celle du \omterm{def:g9:inside-the-atom:nucleus}{noyau}.

\textbf{12.} $3.0 \times 10^{-7} / 9.35 \times 10^{-26} \approx 3.2
\times 10^{18}$ \omterm{def:g9:ions:ion}{ions} fer par litre.
\end{solution}
